\section{Experimental Results}
\label{sec:exp}

\subsection{Segmentation: Quantitative Results}

\paragraph{Per-class mIoU.}
Table~\ref{tab:seg_results} summarizes the segmentation performance of the fine-tuned SegFormer-B0 on the FoodSeg103 validation split.
The 103-class mean IoU is $14.7\%$, which is low in absolute terms but consistent with the extreme long tail of the benchmark: many classes (e.g.\ rare nuts, exotic fruits) appear in fewer than $50$ training images, and per-class IoU on unseen or underrepresented classes is $0$ or NaN by construction.
The top-performing individual classes are broccoli ($80.1\%$), carrot ($71.7\%$), and corn ($68.9\%$), all of which are visually distinct and frequently occurring.

\begin{table}[h]
\centering
\caption{SegFormer-B0 validation results on FoodSeg103 after 5 training epochs.}
\label{tab:seg_results}
\begin{tabular}{lc}
\toprule
Metric & Value \\
\midrule
mIoU (103 classes) & 14.7\% \\
mIoU (7 food groups) & 51.3\% \\
Top class IoU (broccoli) & 80.1\% \\
\bottomrule
\end{tabular}
\end{table}

\paragraph{Per-group mIoU.}
Because SmartPlate's health score depends on group-level proportions rather than individual class labels, the 7-group mIoU is the more meaningful metric.
Table~\ref{tab:group_iou} shows per-group IoU obtained by remapping both predictions and ground-truth through the class-to-group lookup before computing the confusion matrix.
The 7-group mIoU of $51.3\%$ indicates that the model's group-level segmentation is substantially stronger than its fine-grained performance, as expected given that group-level categories have more training support and larger coherent regions.
Vegetables ($67.8\%$) and whole grains ($66.1\%$) score highest; healthy protein ($32.4\%$) is the weakest group, reflecting the visual diversity of protein sources (nuts, eggs, fish, tofu, poultry) and the relatively modest size of their individual pixel regions.

\begin{table}[h]
\centering
\caption{Per-food-group IoU on the FoodSeg103 validation split (SegFormer-B0).}
\label{tab:group_iou}
\begin{tabular}{lc}
\toprule
Food Group & IoU \\
\midrule
Vegetables          & 0.678 \\
Whole Grains        & 0.661 \\
Fruits              & 0.557 \\
Red/Processed Meat  & 0.504 \\
Refined Grains      & 0.450 \\
Sugary/Fatty        & 0.415 \\
Healthy Protein     & 0.324 \\
\midrule
\textbf{Mean}       & \textbf{0.513} \\
\bottomrule
\end{tabular}
\end{table}

\subsection{Baseline Comparison: ResNet-50 vs.\ SegFormer}

The ResNet-50 multi-label classifier achieves a mean F1 of $0.784$ across the 7 food groups after 10 epochs of training, indicating that image-level features are sufficient to detect \emph{which} groups are present on a plate with reasonable accuracy.

However, this result highlights the fundamental limitation of classification-based approaches for dietary assessment.
Table~\ref{tab:comparison} contrasts what each model provides on a representative example.
ResNet-50 correctly identifies that a plate contains vegetables and protein, but reports only binary presence; it misses the fruit component and provides no proportion information.
By contrast, the SegFormer produces continuous proportions---$68.8\%$ vegetables, $7.8\%$ healthy protein, $18.0\%$ refined grains, $5.4\%$ fruits---from which a health score of $56/100$ can be computed directly.
A plate with $70\%$ unhealthy food that ResNet-50 reports as ``vegetables present'' would score misleadingly high under any classification-derived rubric; only pixel proportions surface the imbalance.

\begin{table}[h]
\centering
\caption{What each model provides for a representative validation plate.}
\label{tab:comparison}
\begin{tabular}{lccc}
\toprule
Food Group & ResNet-50 & SegFormer (proportion) \\
\midrule
Vegetables         & present & 68.8\% \\
Fruits             & absent  & 5.4\%  \\
Whole Grains       & absent  & ---    \\
Healthy Protein    & present & 7.8\%  \\
Refined Grains     & present & 18.0\% \\
Red/Processed Meat & absent  & ---    \\
Sugary/Fatty       & absent  & ---    \\
\midrule
Health Score       & \emph{not computable} & 56/100 \\
\bottomrule
\end{tabular}
\end{table}

\subsection{Score Distribution}

Running the full SmartPlate pipeline over a random sample of $200$ held-out validation images yields a mean health score of approximately $48$ with a standard deviation of $18$, reflecting the broad range of plate compositions in FoodSeg103.
Plates dominated by vegetable-heavy Asian dishes score above $70$, while carbohydrate-heavy or dessert plates typically fall below $30$.
The score distribution is roughly unimodal and centered slightly below the midpoint, consistent with FoodSeg103 containing a mix of healthy and indulgent real-world meals rather than any curated healthy eating subset.

\subsection{Qualitative Results}

Figure~\ref{fig:qualitative} shows representative validation plates with the predicted segmentation overlay, group breakdown bar chart, and health score.
The model correctly segments large contiguous food regions (rice, broccoli, salad) but struggles with small or visually ambiguous components such as garnishes, sauces pooling around a protein, or mixed stir-fry dishes where multiple ingredients overlap in the same region.
The health score responds intuitively to the visible plate content: a stir-fried vegetable plate scores $73/100$, while a pizza-heavy plate scores $21/100$ due to the large refined-grain and sugary/fatty penalty.
